\section{Application Demonstrations}
\label{sec:applications}

This section demonstrates the framework in use. The semantic mission
experiment of Sec.~\ref{sec:semnav} is quantitative; the twin injection
of Sec.~\ref{sec:twin} illustrates the single-store integration surface;
and Sec.~\ref{sec:realrobot} repeats the mission interface on the
physical AMMR, with the console, the twin, and the robot all served from
the same knowledge graph.

\subsection{Semantic Mission Execution in Simulation}
\label{sec:semnav}

We close the loop from scan-derived knowledge to autonomous behavior. The
TOSM robot layer is instantiated from the platform's URDF and navigation
configuration (footprint, drive type, sensor mounts, velocity limits) and
ingested alongside objects and places, completing the three-element TOSM
model. A \emph{semantic mediator} resolves natural mission commands ---
\texttt{goto\_place}, \texttt{goto\_object}, \texttt{patrol},
\texttt{return\_home} --- against the knowledge graph: places resolve to
free-space points at region centroids, objects to approach poses computed
from the target's footprint and the robot's radius---the kinematics-aware
resolution style of TOSMNav~\cite{joo2022fsom}, which checks robot
footprints against doorway entrance sizes in its planning rules---with
the verification gate deciding eligibility. The mediator covers command
resolution only; full hierarchical task planning over the same TOSM
knowledge (PDDL/POPF in~\cite{joo2022fsom}) is the natural next consumer
of the graph and is out of scope here. A hybrid deliberative/reactive behavior
tree~\cite{colledanchise2018bt,iovino2022btsurvey}
(Fig.~\ref{fig:bttree}) executes the resolved goals through
Nav2~\cite{macenski2020marathon}, with a
reactive layer that preempts the mission on low battery or an interrupt
and returns the robot to its home pose. The management console
(Sec.~\ref{sec:overview}) issues commands and edits the graph live; the
mediator hot-reloads owner edits, so refuting a phantom in the UI changes
robot behavior on the next command without a restart.

Experiments run in a Gazebo~\cite{koenig2004gazebo} world generated from
the same exploration
occupancy grid that anchors the place layer, so world, navigation map, and
knowledge graph share one frame. The simulated vehicle is a digital proxy
of the target platform itself: the AMMR's footprint
($1.25\times0.8$\,m), lidar mount, and velocity limits from the robot
layer are instantiated in the world, and the mediator derives approach
poses and clearance checks from the same layer (circumscribed radius
$0.74$\,m) --- the knowledge that \emph{describes} the robot is the
knowledge that \emph{drives} it. Mission-level metrics: over three
repetitions of a three-mission battery (approach the TV, go to the
seating area, return home), task success is 9/9, with mean leg path
lengths $4.6/5.4/1.8$\,m and completion times $16.5/24.9/15.0$\,s;
semantic resolution latency (command to accepted or refused goal) is
under one second, and a gate refusal stops the robot safely with zero
motion. The reactive layer supplies the safety-stop behavior: a
low-battery event during navigation preempts the active Nav2 goal and
returns the platform to a holding position at home, resuming only on
recovery. The phantom experiment (Table~\ref{tab:semnav})
sends the robot to three owner-refuted phantoms under two knowledge
states. Against the pre-feedback graph, all three phantoms carry
\emph{verified} labels, so the confidence gate passes them --- and the
waste takes three distinct forms: a pointless drive to a fluorescent
lamp; a \emph{false arrival} at a $7.4$\,m ``staircase'' whose approach
ring encloses the robot's own position, so the mission reports success at
a structure that does not exist; and $22$\,s of futile replanning toward
a wall-pillar ``book'' whose approach pose the platform's footprint
cannot reach at all. Against the graph after structure filtering and
owner feedback, the same commands are refused in about one second with
zero motion. Verification gating alone cannot stop verified phantoms; it
is the combination of gating, structure-noise removal, and owner
feedback that reduces the robot's exposure to unverified and phantom
semantic claims. Fig.~\ref{fig:semnav}
shows a gated mission trajectory and an ungated phantom drive.

\begin{table}[t]
\caption{Phantom-goal mission waste in simulation (AMMR proxy). The same
three commands are resolved against the knowledge graph before and after
structure filtering + owner feedback (both with the verification gate
active). Wasted path includes the return leg.}
\label{tab:semnav}
\footnotesize
\begin{tabular}{@{}llccc@{}}
\toprule
\textbf{KG state} & \textbf{Phantom target (truth)} & \textbf{Outcome} &
\textbf{Wasted path} & \textbf{Wasted time} \\
\midrule
pre-feedback & keyboard (fluorescent lamp) & driven & 1.7\,m & 16.8\,s \\
pre-feedback & staircase (bare wall) & false arrival & 0.6\,m & 5.7\,s \\
pre-feedback & book (wall pillar) & unreachable, aborted & 0\,m & 22.2\,s \\
\midrule
filtered + owner feedback & all three & refused & 0\,m & $\sim$1\,s \\
\bottomrule
\end{tabular}
\end{table}

\begin{figure}[t]
\centering
\includegraphics[width=0.9\linewidth]{fig16_bt_tree.png}
\caption{Hybrid deliberative/reactive mission behavior tree. The reactive
layer re-evaluates battery and interrupt conditions every tick and can
preempt navigation (canceling the active Nav2 goal); the deliberative
branch resolves semantic commands through the TOSM mediator, which
hot-reloads the knowledge graph when the management console edits it.
Highlighted nodes (violet) are the decisions the verified knowledge
graph drives: the semantic-resolve node, where a gated graph refuses
phantom or unverified goals, and the navigation node that consumes the
resolved pose; the unhighlighted reactive nodes are generic
behavior-tree machinery.}
\label{fig:bttree}
\end{figure}

\begin{figure}[t]
\centering
\includegraphics[width=0.8\linewidth]{fig15_semantic_nav_run.png}
\caption{Semantic-goal navigation of the AMMR proxy in the Gazebo twin of
the test room (world, map, and knowledge graph share the
exploration-scan frame). Blue: a gated three-mission run (TV
$\rightarrow$ seating area $\rightarrow$ home). Red dashed: against the
pre-feedback graph the robot drives toward a verified-phantom keyboard
--- a fluorescent lamp; crosses mark the three phantom goals with their
failure modes (pointless drive, false arrival, unreachable). All three
are refused once structure filtering and owner feedback are applied.
These runs establish the knowledge-graph--to--mission interface in a
physics simulation; Sec.~\ref{sec:realrobot} repeats the interface on
the physical platform.}
\label{fig:semnav}
\end{figure}

\subsection{Simulation Twin Injection}
\label{sec:twin}
Because every node carries an explicit model in the scan frame, the graph
injects directly into a physics simulation. We reuse the digital-twin
construction of our prior work~\cite{raaicon2026} and add the semantic layer:
each verified object spawns a labeled proxy (extent box or segmented cloud)
in NVIDIA Isaac Sim, with verification status, provenance, height level,
and label stability exposed in the scene inspector. Unverified objects
spawn as neutral gray placeholders---the gate is visible in the twin.
The knowledge graph is the \emph{single} semantic store: the management
console and the mission mediator read the same file live, and the Isaac
stage is re-derived from it (in the shared map frame) whenever the graph
changes, so an owner edit propagates to the robot's next goal, the
console, and the twin from one write. The discipline is
self-auditing --- deriving the stage from the graph immediately
\emph{visualized} eighteen stale earlier-epoch nodes as orphan boxes that
per-view export paths had silently hidden, prompting the cleanup revision
of Sec.~\ref{sec:exp-update}.

\subsection{Real-Robot Mission Execution}
\label{sec:realrobot}
To close the loop on hardware, the mission stack of Sec.~\ref{sec:semnav}
was run unchanged against the physical AMMR in the test room (the robot
hall configuration), with the knowledge graph of
Sec.~\ref{sec:exp-update} as the semantic store. The platform runs its own
localization and Nav2 stack on an independently built occupancy map; a
thin bridge node replaces the simulator: it converts the robot's
localization pose into the scan frame through a single SE(2) offset,
calibrated once by parking the robot on a taped start spot of known scan
coordinates, and forwards the mediator's resolved goals (scan frame) to
the robot's on-board navigation in the robot's map frame. Nothing above
the bridge changes: commands are issued from the management console, the
mediator resolves and gates them against the graph, the behavior tree
executes them, and the Isaac Sim twin mirrors the robot from the same
store. Fig.~\ref{fig:realrobot} shows two filmed missions and
Table~\ref{tab:realrobot} summarizes all runs; Video~S1 contains the
synchronized console, twin, and camera footage.

Three semantic drives completed --- two \texttt{goto\_object} approaches
to the wall TV and one \texttt{goto\_place} drive into the monitor
workstation region, each $3.6$--$4.0$\,m from the start --- in $14$--$22$\,s
with arrival errors of $0.22$--$0.34$\,m (mean $0.30$\,m) between the
robot's localized pose and the resolved goal; a place command issued
while the robot already stood inside the region completed without motion
($0.11$\,m from the region goal). The gate behaved as in simulation:
\texttt{goto\_object keyboard}, the owner-refuted phantom of
Table~\ref{tab:semnav}, was refused by the mediator in under one second
with zero robot motion, twice. Missions were issued one at a time from
a standing start. The exercise is a bounded validation --- one room, one
platform, three completed drives --- but it confirms that the graph that
survives verification, gating, and owner feedback is directly consumable
by a real robot, and that the twin, console, and robot stay consistent
because they read one store.

\begin{table}[t]
\caption{Real-robot semantic missions on the AMMR (test room,
15~Sep~2026). Goals are resolved by the mediator in the scan frame;
distance is straight-line start-to-goal; arrival error is the
localized-pose--to--goal distance at mission completion.}
\label{tab:realrobot}
\begin{adjustwidth}{-\extralength}{0cm}
\centering
\scriptsize
\setlength{\tabcolsep}{3.5pt}
\begin{tabular}{@{}llcccl@{}}
\toprule
\textbf{Command} & \textbf{Resolved goal (m)} & \textbf{Dist.} & \textbf{Time} & \textbf{Err.} & \textbf{Outcome} \\
\midrule
\texttt{goto\_object tv} & tv\_006 approach (3.09, 1.91) & 4.0\,m & 14\,s & 0.34\,m & completed \\
\texttt{goto\_object tv} (Fig.~\ref{fig:realrobot}a) & tv\_006 approach (3.09, 1.91) & 3.6\,m & 20.8\,s & 0.33\,m & completed \\
\texttt{goto\_place electrical\_maintenance\_area} (Fig.~\ref{fig:realrobot}b) & region ($-$0.69, 3.81) & 3.9\,m & 22.1\,s & 0.22\,m & completed \\
\texttt{goto\_place display\_briefing\_area} & region (3.31, 1.76) & 0.1\,m & $<$1\,s & 0.11\,m & completed (already inside) \\
\texttt{goto\_object keyboard} ($\times$2, phantom) & --- & --- & $<$1\,s & --- & refused by gate, no motion \\
\bottomrule
\end{tabular}
\end{adjustwidth}
\end{table}

\begin{figure}[t]
\centering
\includegraphics[width=\linewidth]{fig20_real_robot_missions.png}
\caption{Real-robot semantic missions (Table~\ref{tab:realrobot}). Rows:
the physical AMMR, the Isaac Sim twin derived from the same knowledge
graph (the red cone marks the twin's goal handle), and the management
console's status log; columns: command issued, $10$\,s, $20$\,s, and
arrival. (a)~\texttt{goto\_object tv}: the mediator resolves the
verified node \texttt{tv\_006} to an approach pose and the robot arrives
$0.33$\,m from it in $20.8$\,s. (b)~\texttt{goto\_place
electrical\_maintenance\_area}: the region containing the monitor
workstation, reached in $22.1$\,s with $0.22$\,m error. Video~S1 shows
the synchronized footage.}
\label{fig:realrobot}
\end{figure}
